\section*{Appendix B: Where everything lives}
\phantomsection
\addcontentsline{toc}{section}{Appendix B: Where everything lives}
\label{app:where}

Every number in this report traces to a devlog entry and an artifact. Code and documentation paths are relative to the repository root (\texttt{sc\_classification/}). Artifact paths are relative to the experiment directory EXP (\nolinkurl{experiments/20260525_060508_stage0_mrd_old34_broad_screen_82db5093/}, about 199 GB, gitignored, on the lab server). Devlog entries are \nolinkurl{scripts/mrd_stage0_2/research/devlog/2026/<id>.md}; \texttt{python tools/research\_index.py show <ID>} prints any devlog, analysis (D/F), hypothesis (H-), task (W\#.n) or ask (A-) ID with everything that cites it.

\subsection*{Section by section}

In Table~\ref{tab:where}, code is under \nolinkurl{scripts/mrd_stage0_2/}, artifact directories are under \nolinkurl{EXP/analysis/scorecards/stage2_diagnostic/}, and notebooks are under \nolinkurl{scripts/mrd_stage0_2/notebooks/stage0_2/}. Each section's figures are regenerated by \texttt{figs\_NN\_<topic>.py} in this report's directory (section 2b: \texttt{figs\_02b\_d26.py}).

{\footnotesize
\begin{longtable}{>{\raggedright\arraybackslash}p{2.6cm} >{\raggedright\arraybackslash}p{4.0cm} >{\raggedright\arraybackslash}p{8.4cm}}
\caption{Where each section's evidence lives.}\label{tab:where}\\
\toprule
Section & Devlog entries & Code, artifact dirs, notebook \\
\midrule
\endfirsthead
\toprule
Section & Devlog entries & Code, artifact dirs, notebook \\
\midrule
\endhead
\bottomrule
\endfoot
1. Cohort and yardstick (D6, D8) & 2026-08-31\_d6-d9-meeting-followup & \emph{Code:} \nolinkurl{stage2_supervised/run_stage2_mrd_multiobjective_scorecard.py}\newline \emph{Artifacts:} \nolinkurl{meeting_followup_v1/}\newline \emph{Notebook:} \nolinkurl{stage2_old34_meeting_followup_metric_interactions.ipynb} \\
\addlinespace
2. Do the highest-prevalence patients shape the factors? (D16, D17, D28) & 2026-10-08\_d16-factor-patient-matrix; 2026-10-08\_d17-basis-variant-pilot; 2026-10-08\_d17-composition-matched-controls; 2026-10-09\_d28-celltype-score-shares & \emph{Code:} \nolinkurl{workstreams/W1_patient_aggregation/run_d16_factor_patient_matrix.py}; \nolinkurl{workstreams/W2_stage1_basis_variants/run_d17_basis_variant_pilot.py}; \nolinkurl{workstreams/W1_patient_aggregation/run_d28_celltype_score_shares.py}\newline \emph{Artifacts:} \nolinkurl{factor_patient_matrix_v1/}; \nolinkurl{basis_variant_pilot_ifna_v1/}; \nolinkurl{factor_patient_matrix_celltype_v1/}\newline \emph{Notebook:} -- \\
\addlinespace
2b. What did we do with allograft rejection and apoptosis? (D26) & 2026-10-09\_d26-leverage-allograft-apoptosis (run entry, decision rule); 2026-10-10\_d26-leverage-allograft-apoptosis-done (results, verdict) & \emph{Code:} \nolinkurl{workstreams/W2_stage1_basis_variants/run_d26_leverage_panels.py} (and .sh for screen \texttt{W2\_d26}); this report's \nolinkurl{figs_02b_d26.py}\newline \emph{Artifacts:} \nolinkurl{leverage_allograft_apoptosis_v1/} (scorecards); the run's own figures under \nolinkurl{EXP/analysis/figures/stage2_diagnostic/leverage_allograft_apoptosis_v1/}\newline \emph{Notebook:} -- \\
\addlinespace
3. Should we aggregate patients? (D14) & 2026-10-02\_d14-patient-factor-overlap & \emph{Code:} \nolinkurl{stage2_supervised/patient_factor_overlap.py}\newline \emph{Artifacts:} \nolinkurl{patient_factor_overlap_v1/}\newline \emph{Notebook:} \nolinkurl{stage2_old34_patient_factor_overlap.ipynb} \\
\addlinespace
4. IFN-$\alpha$: why FA40 reaches only P03 and P04 (D19, D17, D16, F3, F6, F9) & 2026-10-08\_d19-regularization-contract; 2026-10-08\_d16-factor-patient-matrix; 2026-09-17\_team-feedback-f1-f7; 2026-10-02\_f9-ifna-pc8-loadings & \emph{Code:} \nolinkurl{workstreams/W3_dr_value_regularization/regularization_contract.py}; \nolinkurl{workstreams/W3_dr_value_regularization/run_d19_regularization_contract.py}\newline \emph{Artifacts:} \nolinkurl{regularization_contract_v1/}; \nolinkurl{factor_patient_matrix_v1/}; \nolinkurl{feedback_deep_dive_20260917/}\newline \emph{Notebook:} \nolinkurl{stage2_old34_feedback_deep_dives_20260917.ipynb} \\
\addlinespace
5. What is PC8? (D18, D27, F9) & 2026-10-08\_d18-factor-cards; 2026-10-09\_d27-factor-modality-assoc; 2026-10-02\_f9-ifna-pc8-loadings; 2026-10-08\_w7-atlas-strip-factor-explorer & \emph{Code:} \nolinkurl{workstreams/W5_factor_interpretation/run_d18_factor_cards.py}; \nolinkurl{workstreams/W5_factor_interpretation/run_d27_factor_modality_assoc.py} (library \nolinkurl{factor_modality_assoc.py}); \nolinkurl{apps/old34_factor_explorer.py}\newline \emph{Artifacts:} \nolinkurl{factor_cards_v1/}; \nolinkurl{factor_modality_assoc_v1/}; \nolinkurl{feedback_deep_dive_20260917/}\newline \emph{Notebook:} \nolinkurl{stage2_old34_feedback_deep_dives_20260917.ipynb} (section 9) \\
\addlinespace
6. Option 2: inductive Stage 1 (D17) & 2026-10-08\_d17-basis-variant-pilot & \emph{Code:} \nolinkurl{workstreams/W2_stage1_basis_variants/run_d17_basis_variant_pilot.py}; \nolinkurl{workstreams/W2_stage1_basis_variants/stage1_refit.py}\newline \emph{Artifacts:} \nolinkurl{basis_variant_pilot_ifna_v1/}\newline \emph{Notebook:} -- \\
\addlinespace
7. Learning curves (D15, D15b) & 2026-10-02\_d15-cell-count-learning-curve; 2026-10-08\_d15b-patient-count-curve; 2026-10-08\_d15b-full-design & \emph{Code:} \nolinkurl{stage2_supervised/cell_count_learning_curve.py}; \nolinkurl{workstreams/W6_learning_curves/run_d15b_patient_count_curve.py} (and .sh for screen)\newline \emph{Artifacts:} \nolinkurl{cell_count_learning_curve_v1/}; \nolinkurl{patient_count_learning_curve_v1/}\newline \emph{Notebook:} \nolinkurl{stage2_old34_patient_factor_overlap.ipynb} (D15) \\
\addlinespace
8. The three scientific questions (D1, D1e, D1f, D3, D6, D9, D12--D19, F2/F4) & 2026-09-11\_d1-atlas-783-d1e-d1f; 2026-08-31\_d6-d9-meeting-followup; 2026-09-10\_d1a-descriptive-deep-dives; 2026-09-17\_team-feedback-f1-f7; 2026-09-03\_d12-d13-gsea-hvg-lowsupport & \emph{Code:} \nolinkurl{stage2_supervised/stage2_sharedness_plotting.py}\newline \emph{Artifacts:} the directory itself (\nolinkurl{old34_*} tables); \nolinkurl{meeting_followup_v1/}; \nolinkurl{d1a_deep_dive_20260910/}; \nolinkurl{deg_hvg_lowsupport_v1/}\newline \emph{Notebook:} \nolinkurl{stage2_old34_comprehensive_diagnostic_suite.ipynb}; \nolinkurl{stage2_old34_meeting_followup_metric_interactions.ipynb} \\
\addlinespace
9. Remaining TODOs and decisions (D1DD13, D6, D15b) & 2026-09-10\_d1a-descriptive-deep-dives; 2026-10-08\_d15b-full-design; 2026-10-08\_literature-reviews; 2026-10-08\_team-update-and-wave2-integration; 2026-10-09\_report-reading-notes; 2026-10-09\_w8-healthy-downsampling-plan & \emph{Code:} \nolinkurl{workstreams/W8_new_cohort/} (email draft, integration plan with the patient-inclusion section, healthy-cell downsampling plan); \nolinkurl{workstreams/W9_literature/} (reviews)\newline \emph{Artifacts:} \nolinkurl{d1a_deep_dive_20260910/}; \nolinkurl{meeting_followup_v1/}; \nolinkurl{patient_count_learning_curve_v1/}\newline \emph{Notebook:} \nolinkurl{stage2_old34_d1a_descriptive_deep_dives_20260910.ipynb} \\
\end{longtable}
}

\subsection*{The project's context layer}

\begin{tabularx}{\linewidth}{>{\raggedright\arraybackslash}p{5.6cm} >{\raggedright\arraybackslash}X}
\toprule
File & What it is \\
\midrule
\texttt{AGENTS.md} & Start here: the repository in brief and the hard rules R1--R12 (long jobs only in your screen, retune C whenever training $n$ changes, claim language) \\
\nolinkurl{scripts/mrd_stage0_2/research/STATUS.md} & Generated dashboard: cohort, workstreams, open tasks by priority, hypotheses and current verdicts, analysis registry, campaigns \\
\texttt{research/THEME.md} & The central question, Q1--Q3, and the canonical claim boundary (\S6) \\
\texttt{research/CLAIMS.md}, \texttt{hypotheses.md} & Hypotheses with their evidence and verdicts \\
\texttt{research/devlog/} & One entry per result: question, estimand, code, artifacts, Results, claim gate, next steps \\
\texttt{research/analyses.csv}, \texttt{campaigns.csv} & Analysis registry (ID $\to$ code, notebook, artifact dir) and compute campaigns with their resume commands \\
\texttt{research/asks/}, \texttt{team\_updates/} & Questions as asked and dated team updates. This report answers A-20261008-01 to -32 (\texttt{asks/\allowbreak 2026-10-08\_team-report.md}); your reading notes on version 1 are A-20261009-01 to -16 (\texttt{asks/\allowbreak 2026-10-09\_report-reading-notes.md}), with a table of where version 2 takes each up \\
\texttt{research/paper/} & The living paper context, drafted alongside the runs: \texttt{README.md} (rules), \texttt{OUTLINE.md} (working titles, results arcs mapped to hypotheses, analyses and devlog ids, methods, experiments still needed) and \texttt{FIGURES.md} (the panel plan). Number-free; THEME \S6 stays the claim boundary \\
\nolinkurl{scripts/mrd_stage0_2/workstreams/} (W0--W9) & Each workstream's charter, backlog (W\#.n tasks), code and conventions \\
\bottomrule
\end{tabularx}

\subsection*{Regenerating this report}

From the repository root, in the lab env (\nolinkurl{sc-classification-rapids-py311}) with the experiment directory available:
\begin{quote}\raggedright
\texttt{bash} \nolinkurl{scripts/mrd_stage0_2/research/reports/2026-10-08_questions_report/build.sh} \texttt{-{}-full}
\end{quote}
This writes \texttt{report\_full.pdf}; without \texttt{-{}-full} the script builds the short core report, \texttt{report\_core.pdf}, and \texttt{-{}-all} builds both plus the team version (\texttt{README.md} in the report directory).
Each figure module can also be run alone, for example \texttt{python} \nolinkurl{scripts/mrd_stage0_2/research/reports/2026-10-08_questions_report/figs_09_timing.py}: its \texttt{main()} rewrites that section's \texttt{figures/<name>.pdf}, plus a \texttt{.png} preview, from read-only tables in EXP. The modules share \texttt{report\_style.py} (palette, type sizes, output paths) and resolve EXP through \nolinkurl{sc_classification.utils.paths}, so the report directory holds no absolute paths. Sections are self-contained: to share a subset, remove the \texttt{\textbackslash input} lines you do not want from \texttt{report.tex} and rebuild. \texttt{build.sh -{}-team} writes \texttt{report\_team.pdf}, which leaves out Appendix~C (your repository decisions).

\subsection*{Browsing the results interactively}

Two read-only Dash apps live in \nolinkurl{scripts/mrd_stage0_2/apps/}. Both bind to 127.0.0.1 by default; to share from the lab server, forward the ports over SSH (for example \texttt{ssh -L 8050:localhost:8050 -L 8051:localhost:8051 <server>}).
\begin{itemize}
\item \textbf{D1 atlas explorer}, \texttt{old34\_d1\_atlas\_explorer.py} (port 8050): the 47 Stage 0 panels $\times$ 17 Stage 1 configurations $\times$ 7 supported patients as patient-specific heatmaps of lift or threshold counts. A cell's view shows the per-fold strip, with all-zero folds marked ``lift = 0 by construction'', the selected penalty and C, and an ``Open factor cards'' link.
\item \textbf{Factor explorer}, \texttt{old34\_factor\_explorer.py} (port 8051): the 1,810 D18 factor cards over 56 representations, with both gene poles, per-patient univariate lift, the fold coefficient strip, and ADT and cell-type tables. URLs address a card directly, for example \nolinkurl{?panel=single_geneset__hallmark_interferon_alpha_response&method=pca&k=40&factor=8} for PC8.
\item \textbf{Checks:} \texttt{-{}-check} validates either app's inputs against the stored tables and prints a pass/fail report without serving.
\item \textbf{Without EXP:} \texttt{-{}-export-bundle DIR} writes a small shareable copy to a new directory (about 2.2 MB of tables for the atlas, about 7.0 MB for the factor cards); a teammate runs the same script with \texttt{-{}-data-dir DIR}. For everything else, \texttt{STATUS.md}, \texttt{CLAIMS.md} and the devlog are in git.
\item \textbf{Long-running servers} go in your screen sessions (\texttt{d1\_atlas} and \texttt{W7\_factor\_explorer}), as recorded in devlog 2026-10-08\_w7-atlas-strip-factor-explorer.
\end{itemize}
